\documentclass[11pt]{article}

\usepackage[a4paper,margin=1in]{geometry}
\usepackage{amsmath,amssymb,amsthm,mathtools}
\usepackage{booktabs}
\usepackage[hidelinks]{hyperref}

\newtheorem{theorem}{Theorem}[section]
\newtheorem{proposition}[theorem]{Proposition}
\newtheorem{lemma}[theorem]{Lemma}
\newtheorem{corollary}[theorem]{Corollary}
\theoremstyle{definition}
\newtheorem{definition}[theorem]{Definition}
\newtheorem{remark}[theorem]{Remark}

\newcommand{\R}{\mathbb{R}}
\newcommand{\Fr}{\mathcal{A}}
\newcommand{\Om}{\Omega}
\newcommand{\Otilde}{\widetilde{\Omega}}
\newcommand{\eps}{\varepsilon}
\newcommand{\rstar}{\rho_{*}}
\newcommand{\rhad}{\widehat{\rho}}
\newcommand{\zbar}{\bar z_{\rho}}
\newcommand{\X}{\mathcal{X}}
\newcommand{\Bstar}{B^{*}}
\newcommand{\PiE}{\Pi(E_{\rho},\zbar)}

\title{Sharp Quantitative Faber--Krahn Stability via Route~$\delta$\\
       (Plan~2: a No-Selection, No-Graph, No-Schauder Proof)}
\author{Plan~2 / Route~$\delta$: assembled document}
\date{}

\begin{document}
\maketitle

\begin{abstract}
We give a proof of the sharp quantitative Faber--Krahn inequality
\[
   \Bigl(\frac{|\Om|}{|B_1|}\Bigr)^{2/n}
   \bigl(\lambda_1(\Om)-\lambda_1(\Bstar)\bigr)
   \;\ge\;c_{\rm FK}(n)\,\Fr(\Om)^2,
\]
with sharp exponent~$2$ on the Fraenkel asymmetry $\Fr$
(equivalently $1/2$ on the scale-normalised spectral deficit) that
is structurally distinct from the
Brasco--De~Philippis--Velichkov~\cite{BraDeP2017} proof and from the
contradiction-free Plan~1 / Route~A proof~\cite{Plan1Master}. The
present proof, which we call \emph{Route~$\delta$}, makes no use of
any of the following four classical ingredients: (i) the
selection / penalisation principle of Cicalese--Leonardi and BDV;
(ii) the graph-entry / Hausdorff-to-$C^{1,\gamma}$ step driven by
free-boundary regularity; (iii) Schauder bootstrap with
Gagliardo--Nirenberg interpolation; (iv) the nearly-spherical
Fuglede-type closure theorem. In their place it uses a
finite-perimeter level-set deficit identity, a metric foliation
$\rho\mapsto[F_\rho]\in\X$ in the $L^1$-quotient $\X$ of
characteristic classes modulo translations, a centroid kinematic
identity, a space--time radial-moment identity, and a
slicing-then-squaring decomposition that converts an integrated
quadratic boundary oscillation into a sharp metric trace. The proof
is organised into seven self-contained building blocks
(\textsc{LevelSetIdentity}, \textsc{MetricFramework},
\textsc{CentroidBound}, \textsc{SpaceTimeIdentity},
\textsc{WeightedMetricTrace}, \textsc{BoundaryLayerTransfer},
\textsc{GlobalAssembly}), each delivered as a separate \LaTeX{}
document. The last two stages of the chain (bounded reduction and
the Kohler--Jobin rearrangement) are imported verbatim from the
Plan~1 Final/ library, so that the methodological novelty is
concentrated in the first five blocks. The route avoids the
pointwise quadratic Fusco--Julin oscillation~$(B^2)$ entirely; only
its integrated form is consumed.
\end{abstract}

\tableofcontents

\section{Introduction}\label{sec:intro}

\subsection{The Faber--Krahn inequality}\label{ssec:FK}

Let $\Om\subset\R^n$, $n\ge 2$, be an open set of finite Lebesgue
measure and let
\[
   \lambda_1(\Om)\;=\;\min_{u\in H^1_0(\Om)\setminus\{0\}}
                    \frac{\int_\Om|\nabla u|^2\,dx}{\int_\Om u^2\,dx}
\]
be the first Dirichlet eigenvalue of $-\Delta$ on $\Om$. The
\emph{Faber--Krahn inequality} asserts that among open sets of a
prescribed volume the ball uniquely minimises $\lambda_1$: writing
$\Bstar$ for the ball with $|\Bstar|=|\Om|$,
\begin{equation}\label{eq:FKsharp}
   \lambda_1(\Om)\;\ge\;\lambda_1(\Bstar),
\end{equation}
with equality if and only if $\Om$ is a ball (up to a set of capacity
zero). This was conjectured by Lord Rayleigh in the planar case in
the late nineteenth century and proved in general dimension by
Faber~\cite{Faber1923} and Krahn~\cite{Krahn1925} via Schwarz
symmetrisation and the P\'olya--Szeg\H{o} principle; modern accounts
appear in~\cite{KawohlBook,HenrotBook,BucurButtazzo}.

The natural way to measure the deviation from spherical symmetry is
the \emph{Fraenkel asymmetry}
\begin{equation}\label{eq:Fraenkel}
   \Fr(\Om)\;=\;\inf_{x_0\in\R^n}
   \frac{|\Om\,\Delta\,(\Bstar+x_0)|}{|\Bstar|}\;\in\;[0,2),
\end{equation}
which vanishes precisely on translates of the ball. A quantitative
Faber--Krahn inequality is an estimate of the form
$(|\Om|/|B_1|)^{2/n}(\lambda_1(\Om)-\lambda_1(\Bstar))\ge
c\,\omega(\Fr(\Om))$ for some modulus $\omega$. The ellipsoidal
perturbation family $\Om_\eps=\{|x'|^2+(1+\eps)x_n^2\le 1\}$ shows
$\Fr(\Om_\eps)\simeq\eps$ and
$\lambda_1(\Om_\eps)-\lambda_1(B^*)\simeq\eps^2$, so the sharp
modulus is~$\Fr^2$.

\subsection{Stability: prior work}\label{ssec:prior}

The first systematic study of Faber--Krahn stability is due to
Hansen--Nadirashvili~\cite{HansenNadirashvili1994} and
Melas~\cite{Melas1992} for restricted classes of sets. Bhattacharya
\cite{Bhattacharya2001} reached every planar open set with
exponent~$3$. Fusco--Maggi--Pratelli~\cite{FMPLaplace}, using their
sharp quantitative isoperimetric inequality~\cite{FMP2008} (see also
the mass-transport approach of Figalli--Maggi--Pratelli
\cite{FMP2010}), extended this to all dimensions with exponent~$4$
on $\Fr$.

The sharp conjecture with exponent~$2$ was settled in 2015 by
Brasco--De~Philippis--Velichkov~\cite{BraDeP2017}: there exists
$\sigma(n)>0$ with
$|\Om|^{2/n}\lambda_1(\Om)-|\Bstar|^{2/n}\lambda_1(\Bstar)\ge
\sigma\Fr(\Om)^2$. The BDV strategy has two pillars: a
Kohler--Jobin~\cite{KohlerJobin,Brasco2014} reduction of
Faber--Krahn to the analogous sharp Saint--Venant stability, then a
selection principle in the spirit of
Cicalese--Leonardi~\cite{CL2012} and
Acerbi--Fusco--Morini~\cite{AFM2013} combined with free-boundary
regularity~\cite{AC1981} and a nearly-spherical expansion
(classical for the isoperimetric deficit \cite{Fuglede}, extended
to first-eigenvalue deficit in \cite[Thm.~3.3]{BraDeP2017}) to produce
a nearly spherical contradiction. Subsequent work of
Allen--Kriventsov--Neumayer~\cite{AKN} addresses related
spectral stability inequalities via different routes.

A constructive-constant version of the BDV theorem was given in
Plan~1~\cite{Plan1Master} (\emph{Route~A}): the BDV
selection / contradiction is replaced by a quantitative
\emph{single-set} selection (run once on a single near-extremal
set), the regularity package is applied to that single set, and
the BDV nearly spherical theorem is invoked directly; no
subsequence is taken. Plan~1 retains, however, all four of the BDV
technical ingredients (selection, graph entry, Schauder,
nearly-spherical closure); it removes only the contradiction.

The persistent open problem the present paper addresses is whether
a proof of the sharp inequality can be given that uses
\emph{none} of the four ingredients above. The price for this is a
new \emph{finite-perimeter} infrastructure for the level-set
foliation of the torsion function, which is the content of
Plan~2 / Route~$\delta$.

\subsection{Plan~2's contribution}\label{ssec:plan2}

Plan~2 gives a proof of the sharp quantitative Faber--Krahn
inequality in which the small-asymmetry chain
(S)$\to$(G)$\to$(I)$\to$(N) of Plan~1 is replaced by a
\emph{metric level-set chain}. Explicitly:

\begin{enumerate}
\item[\rm(a)] \emph{No selection principle.} No penalised functional
is minimised. The original domain $\Om$ is analysed as it stands;
its volume-radius foliation $E_\rho=\{u_\Om>t(\rho)\}$ is the
structural object.

\item[\rm(b)] \emph{No graph entry.} No claim that $\partial\Om$, or
that any selected $\Otilde$, is a normal graph over a sphere of
class $C^{1,\gamma}$ is made. All boundary integrals are on
\emph{reduced boundaries} $\partial^*E_\rho$ in the De Giorgi
sense.

\item[\rm(c)] \emph{No Schauder interpolation.} No
Kinderlehrer--Nirenberg hodograph straightening, no Lieberman
oblique Schauder bootstrap, no Gagliardo--Nirenberg interpolation.
The proof remains in $W^{2,p}_{\rm loc}$ throughout.

\item[\rm(d)] \emph{No nearly-spherical closure theorem.} BDV
Theorem~3.3 and the Fuglede $H^{1/2}$ expansion are not invoked.
The deficit-to-asymmetry conversion is purely metric.
\end{enumerate}

In their place, the chain uses the following four new ingredients:

\begin{enumerate}
\item[\rm(1)] \emph{The metric foliation.} The torsion volume-radius
parametrisation $E_\rho=\{u>t(\rho)\}$ with $|E_\rho|=\omega_n\rho^n$
defines a curve $\rho\mapsto[F_\rho]$ in the metric space
$(\X,d_\X)$ of $L^1$-classes modulo translations
(Building Block~\textsc{MetricFramework}). The
Ambrosio--Gigli--Savar\'e metric framework~\cite{AGS2008} furnishes
a metric derivative $|\dot F_\rho|_\X$, controlled by a Reshetnyak /
coarea limit of the moving-region first variation.

\item[\rm(2)] \emph{The centroid kinematic identity.} The bulk
centroid $\zbar=|E_\rho|^{-1}\int_{E_\rho}x\,dx$ satisfies
\[
   \bar z_\rho'\;=\;\frac{1}{|E_\rho|}\int_{\partial^*E_\rho}
        (x-\zbar)\,V_\rho\,d\mathcal H^{n-1},
\]
with $V_\rho=-t'(\rho)/|\nabla u|$ the normal velocity. The
right-hand side has \emph{no homothety term}, so the bound on
$\bar z_\rho'$ goes through the unconditional pair
(Cauchy--Schwarz on $V_\rho$, FMP / Fusco--Julin at the centroid)
and does not consume any per-$\rho$ radial-trace input
(Building Block \textsc{CentroidBound}).

\item[\rm(3)] \emph{The space--time $\Pi$ identity.} The signed
radial moment
$\PiE=(n-1)\!\int(\mathbf1_{B_\rho(\zbar)}-\mathbf1_{E_\rho})|x-\zbar|^{-1}dx$
satisfies, by Fubini in $(x,\rho)$,
$\int_{\rstar}^{\rho_\delta}\PiE\,d\mu(\rho)=
(n-1)\int_{\widetilde\Om}\sigma(x)dx$
for an explicit Fubini density~$\sigma$
(Building Block \textsc{SpaceTimeIdentity}).

\item[\rm(4)] \emph{Slicing-then-squaring.} The radial $L^1$ trace
$T_1=\int_{\partial^*E_\rho}|\,|x-\zbar|/\rho-1|\,d\mathcal H^{n-1}$
is split as $T_1=T_1^{\rm rad,slic}+T_1^{\rm tan}$; the radial
piece is bounded \emph{linearly} via the good-ray equality of
slicewise volumes and a bad-ray multiplicity bound; the squared
volume-Fraenkel $|E_\rho\Delta B_\rho|^2$ is controlled by~$\PiE$
via a slice-rearrangement. Squaring before integrating preserves
the sharp $\delta$-rate that the alternative $T_2$-route would
leak at $\sqrt\delta$.
\end{enumerate}

\begin{theorem}[Sharp quantitative Faber--Krahn, Plan~2 form]
\label{thm:mainintro}
Let $n\ge 2$. There exists a constant $c_{\rm FK}(n)>0$, depending
only on the dimension, given explicitly by~\eqref{eq:cFKdef} below
as a composition of named constants of the seven Plan~2 building
blocks (and of the Plan~1 \textsc{BoundedReduction} and
\textsc{KohlerJobinTransfer} blocks at the dimensional radius
$R_*(n)$), such that
\[
   \Bigl(\frac{|\Om|}{|B_1|}\Bigr)^{2/n}
   \bigl(\lambda_1(\Om)-\lambda_1(\Bstar)\bigr)
   \;\ge\;c_{\rm FK}(n)\,\Fr(\Om)^2
\]
for every open $\Om\subset\R^n$ of finite positive measure. The
exponent~$2$ on $\Fr$ is sharp.
\end{theorem}

\subsection{Organisation of the paper}\label{ssec:org}

The argument is organised as a chain of seven self-contained
building blocks. Schematically,
\[
   \textsc{LevelSetIdentity}\to
   \textsc{MetricFramework}\to
   \textsc{CentroidBound}\to
   \textsc{SpaceTimeIdentity}\to
   \textsc{WeightedMetricTrace}\to
   \textsc{BoundaryLayerTransfer}\to
   \textsc{GlobalAssembly}.
\]
The first six blocks are the geometric heart of Plan~2; the seventh
imports the bounded reduction of \cite[\S5]{BraDeP2017} (equivalently
Plan~1's \textsc{BoundedReduction}~\cite{BDV-BoundedReduction-Plan1})
and the Kohler--Jobin transfer of Brasco~\cite{Brasco2014} (Plan~1's
\textsc{KohlerJobinTransfer}~\cite{KohlerJobinTransfer-Plan1}).

\section{Notation and standing hypotheses}\label{sec:notation}

Fix $n\ge 2$ and $R\ge 2$. We use the following normalisations:
\begin{itemize}
\item $u_\Om\in H^1_0(\Om)$ is the variational torsion function:
$-\Delta u_\Om=1$ in $\Om$, $u_\Om=0$ on $\partial\Om$.
\item $T(\Om)=\int_\Om u_\Om$ is the torsional rigidity, and
$E(\Om)=-T(\Om)/2$ is the Saint--Venant energy.
\item $\delta_T(\Om):=E(\Om)-E(\Bstar)\ge 0$ at fixed volume.
\item $\lambda_1(\Om)$ is the first Dirichlet eigenvalue of $-\Delta$.
\item $\Fr(\Om)=\inf_{x_0}|\Om\Delta(\Bstar+x_0)|/|\Bstar|\in[0,2)$.
\item For $\rho\in[\rstar,1]$, $E_\rho=\{u_\Om>t(\rho)\}$ with
$|E_\rho|=\omega_n\rho^n$ and $\mu(d\rho)=(-t'(\rho))d\rho$.
\item $\zbar=|E_\rho|^{-1}\int_{E_\rho}x\,dx$ is the centroid of
$E_\rho$.
\item $F_\rho=\rho^{-1}E_\rho$ is the rescaled level set;
$[F_\rho]\in\X$ its class in $\X=\{L^1\text{-classes}\}/\R^n$ with
$d_\X([A],[C])=\inf_a|A\Delta(C+a)|$.
\item $V_\rho(x)=-t'(\rho)/|\nabla u(x)|$, $H_{z,\rho}(x)=(x-z)\cdot\nu_{E_\rho}(x)/\rho$.
\item $\rho_\delta:=1-\kappa\sqrt{\delta_T(\Om)}$;
$\rhad\in[\rho_\delta-C_*\delta_T,\rho_\delta]$ the metric-trace radius.
\item $\omega_n=|B_1|$, $L_n=\lambda_1(B_1)$, $T_n=T(B_1)=\omega_n/(n(n+2))$.
\item $R_*(n)=\max\{d(n),2\}$ the BDV bounded-reduction radius; we
fix $\rstar=1/2$.
\end{itemize}

\noindent\textbf{Notation convention (two parametrisations).} We use
both notations: $E_t=\{u>t\}$ (level-set parametrisation, Block~I)
and $E_\rho=\{u>t(\rho)\}$ with $|E_\rho|=\omega_n\rho^n$
(volume-radius parametrisation, Blocks~II--VII). The map
$t\leftrightarrow\rho$ is the inverse distribution function.

\section{The seven building blocks chained}\label{sec:chain}

\subsection{Block I: \textsc{LevelSetIdentity}}\label{ssec:L}

\begin{theorem}[Exact deficit identity; \cite{LevelSetIdentity}]
\label{thm:L}
Let $\Om\subset\R^n$ be open with $|\Om|<\infty$. With
$E_t=\{u>t\}$, $m(t)=|E_t|$,
$\alpha(t)=\int_{\partial^*E_t}|\nabla u|^{-1}d\mathcal H^{n-1}$,
$\beta(t)=\int_{\partial^*E_t}|\nabla u|\,d\mathcal H^{n-1}$,
$D_H(t)=\alpha\beta-P(E_t)^2$,
$D_I(t)=P(E_t)^2-n^2\omega_n^{2/n}m(t)^{2-2/n}$,
\begin{equation}\label{eq:Lboxed}
\boxed{\;
   \frac{2}{|\Om|}\bigl(E(\Om)-E(B)\bigr)
   =\frac{1}{|\Om|}\int_{0}^{\|u\|_\infty}
        \frac{D_H(t)+D_I(t)}{n^2\omega_n^{2/n}\,m(t)^{1-2/n}}\,dt.\;}
\end{equation}
Both $D_H,D_I\ge 0$ a.e.; no boundedness or smoothness of $\Om$ is
needed. All integrals on the reduced boundary $\partial^*E_t$.
\end{theorem}

\subsection{Block II: \textsc{MetricFramework}}\label{ssec:M}

\begin{theorem}[Metric first variation; \cite{MetricFramework}]
\label{thm:M}
For $\Om\subset B_R$ with $|\Om|=\omega_n$ and any Borel centre
field $\{z_\rho\}$ with $|z_\rho|\le R'$, there is
$C_{n,\rstar,R}>0$ such that
\begin{equation}\label{eq:Mboxed}
\boxed{\;
   |\dot F_\rho|_\X
   \le C_{n,\rstar,R}\inf_{a\in\R^n}\!\int_{\partial^*E_\rho}
        |V_\rho-H_{z_\rho,\rho}-a\cdot\nu_\rho|\,d\mathcal H^{n-1}
\;}
\end{equation}
for a.e.\ $\rho\in[\rstar,1]$.
\end{theorem}

\subsection{Block III: \textsc{CentroidBound}}\label{ssec:F}

\begin{theorem}[$H^1$-in-$\rho$ centroid bound (F);
\cite{CentroidBound}]\label{thm:F}
$\rho\mapsto\zbar$ is Lipschitz on $[\rstar,1]$ and there is
$C_{\rm cent}(n,R,\rstar)>0$ with
\begin{equation}\label{eq:Fboxed}
\boxed{\;
   \int_{\rstar}^{\rho_\delta}|\bar z_\rho'|^2\,(-t'(\rho))\,d\rho
   \;\le\;C_{\rm cent}\,\delta_T(\Om).\;}\tag{F}
\end{equation}
The proof uses the centroid kinematic identity
$\bar z_\rho'=|E_\rho|^{-1}\int_{\partial^*E_\rho}(x-\zbar)V_\rho\,d\mathcal H^{n-1}$
(no homothety term), the unconditional $D_H$-Cauchy on $V_\rho$,
FMP / Fusco--Julin at the centroid, and the integrated Talenti
gap of Theorem~\ref{thm:L}.
\end{theorem}

\subsection{Block IV: \textsc{SpaceTimeIdentity}}\label{ssec:Pi}

\begin{theorem}[Space-time $\Pi$ identity and integrated bounds;
\cite{SpaceTimeIdentity}]\label{thm:Pi}
Conditional on \eqref{eq:Fboxed}:
\begin{equation}\label{eq:Piboxed}
\boxed{\;
   \int_{\rstar}^{\rho_\delta}\PiE\,d\mu(\rho)
   =(n-1)\!\int_{\widetilde\Om}\sigma(x)\,dx,\;}
\end{equation}
the integrated quadratic oscillation
$\int\!\int|\nu-e_r|^2\,d\mathcal H^{n-1}d\mu\le C_{B^2}\delta_T$
(B$^2$-int), and the integrated radial trace
$\int(\int|H_{\zbar,\rho}-1|d\mathcal H^{n-1})^2\,d\mu\le C_{\rm radial}\delta_T$
(T$_1^2$-int). The proof uses
$T_1=T_1^{\rm rad,slic}+T_1^{\rm tan}$ with linear good-ray
equality and the slice-rearrangement
$|E_\rho\Delta B_\rho|^2\le C\,\PiE$.
\end{theorem}

\subsection{Block V: \textsc{WeightedMetricTrace}}\label{ssec:K}

\begin{theorem}[Weighted metric trace at $\rhad$;
\cite{WeightedMetricTrace}]\label{thm:K}
Conditional on \eqref{eq:Fboxed} and (B$^2$-int),
$\int(d_\X(F_\rho,B_1)^2+|\dot F_\rho|_\X^2)d\mu\le C\,\delta_T$
(integrated action), $\int|\dot F_\rho|_\X^2\,d\rho\le C'\delta_T$
(kinetic bound), and there is
$\rhad\in[\rho_\delta-C_*\delta_T,\rho_\delta]$ with
\begin{equation}\label{eq:Kboxed}
\boxed{\;
   d_\X(F_{\rhad},B_1)^2\;\le\;C_*\,\delta_T(\Om).
\;}
\end{equation}
\end{theorem}

\subsection{Block VI: \textsc{BoundaryLayerTransfer}}\label{ssec:bdyL}

\begin{theorem}[Bounded sharp Saint--Venant;
\cite{BoundaryLayerTransfer}]\label{thm:BL}
There is $C(n,R,\rstar)>0$ such that every open $\Om\subset B_R$
with $|\Om|=\omega_n$ satisfies
\begin{equation}\label{eq:BLboxed}
\boxed{\;
   \Fr(\Om)^2\le C(n,R,\rstar)\,\delta_T(\Om).
\;}
\end{equation}
The proof combines \eqref{eq:Kboxed} with Lemma~8.3 of
\cite{BoundaryLayerTransfer} ($\Fr(\Om)\le\Fr(E_{\rhad})+2|\Om\setminus
E_{\rhad}|/\omega_n$) and the squaring identity
$(a+b)^2\le 2(a^2+b^2)$; sharp exponent~$2$.
\end{theorem}

\subsection{Block VII: \textsc{GlobalAssembly}}\label{ssec:GA}

\begin{theorem}[Sharp Faber--Krahn; \cite{GlobalAssembly}]
\label{thm:Glob}
Combining~\eqref{eq:BLboxed} at the dimensional radius $R_*(n)$ with
the BDV bounded reduction~\cite{BDV-BoundedReduction-Plan1} and the
Kohler--Jobin transfer~\cite{KohlerJobinTransfer-Plan1}, there is a
dimensional constant $c_{\rm FK}(n)>0$ given by
\begin{equation}\label{eq:cFKdef}
   c_{\rm FK}(n):=\min\!\Bigl\{
      \frac{4L_n\,c_{\rm SV}^{\rm glob}(n)}{(n+2)T_n},\,
      \frac{L_n(2^{2/(n+2)}-1)}{4}\Bigr\},
\end{equation}
such that for every open $\Om\subset\R^n$ with $0<|\Om|<\infty$,
\begin{equation}\label{eq:Globboxed}
\boxed{\;
   \Bigl(\frac{|\Om|}{|B_1|}\Bigr)^{2/n}
   \bigl(\lambda_1(\Om)-\lambda_1(\Bstar)\bigr)
   \ge c_{\rm FK}(n)\,\Fr(\Om)^2.
\;}
\end{equation}
\end{theorem}

\section{The main theorem and final assembly}\label{sec:assembled}

\begin{theorem}[Sharp quantitative Faber--Krahn stability,
Plan~2 form]\label{thm:final}
Let $n\ge 2$. There is a dimensional constant $c_{\rm FK}(n)>0$,
given by \eqref{eq:cFKdef}, such that for every open
$\Om\subset\R^n$ of finite positive measure,
$(|\Om|/|B_1|)^{2/n}(\lambda_1(\Om)-\lambda_1(\Bstar))\ge c_{\rm FK}(n)\Fr(\Om)^2$.
The exponent~$2$ on $\Fr$ is sharp.
\end{theorem}

\begin{proof}
By scale invariance of both sides --- $(|\Om|/|B_1|)^{2/n}(\lambda_1(\Om)-\lambda_1(B^*))$
is dimensionless, as is $\Fr(\Om)$ --- WLOG $|\Om|=\omega_n$ after
rescaling $\Om\mapsto(\omega_n/|\Om|)^{1/n}\Om$. Fix $\rstar=1/2$.
The chain Theorems~\ref{thm:L}--\ref{thm:BL} delivers the bounded
Saint--Venant inequality at $R=R_*(n)$.
Theorem~\ref{thm:Glob} then internalises the $R$-dependence via
the BDV bounded reduction and converts to Faber--Krahn via
Kohler--Jobin.
\end{proof}

\section{Constants chain}\label{sec:constants}

\renewcommand{\arraystretch}{1.25}
\begin{center}
\begin{tabular}{|l|l|l|}\hline
Stage & Constant & Dependence \\\hline
(L) deficit identity & --- & --- \\
(M) first variation & $C_{n,\rstar,R}=\rstar^{-n}$ & $n,\rstar,R$ \\
(F) centroid bound & $C_{\rm cent}$ & $n,R,\rstar$ \\
($\Pi$) (B$^2$-int) & $C_{B^2}$ & $n,R,\rstar$ \\
($\Pi$) (T$_1^2$-int) & $C_{\rm radial}$ & $n,R,\rstar$ \\
(K) action, kinetic, trace & $C,C',C_*$ & $n,R,\rstar$ \\
($\partial$) bounded SV & $C(n,R,\rstar)$ & $n,R,\rstar$ \\
$\rstar=1/2$, $R\to R_*(n)$ & $C(n,R_*(n),1/2)$ & $n$ \\
(FK) BDV constants & $C_9(n),\delta(n),d(n)$ & $n$ \\
(FK) global SV & $c_{\rm SV}^{\rm glob}(n)$ & $n$ \\
(FK) KJ multiplier & $4L_n/((n+2)T_n)$ & $n$ \\
(FK) final & $c_{\rm FK}(n)$ & $n$ \\\hline
\end{tabular}
\end{center}

\textbf{Conclusion.} All $R$- and $\rstar$-dependence is absorbed at
the Block~($\partial$) $\to$ (FK) interface; final constant
$c_{\rm FK}(n)$ depends only on $n$.

\section{Comparison with Plan~1 (Route~A)}\label{sec:vsPlan1}

Both routes deliver Theorem~\ref{thm:final}. The structural
difference is concentrated in the bounded sharp Saint--Venant
stage. Plan~1 uses the chain
(S)\,selection $\to$ (G)\,graph entry $\to$ (I)\,Schauder
interpolation $\to$ (N)\,nearly-spherical closure
(\cite{Plan1Master}). Plan~2 uses the chain
(L)\,deficit identity $\to$ (M)\,metric first variation $\to$
(F)\,centroid bound $\to$ ($\Pi$)\,space-time + slicing-then-squaring
$\to$ (K)\,metric trace $\to$ ($\partial$)\,boundary-layer transfer.

\begin{center}
\renewcommand{\arraystretch}{1.18}
\begin{tabular}{|l|c|c|}\hline
Tool & Plan~1 / Route~A & Plan~2 / Route~$\delta$ \\\hline
Selection / penalised functional & uses & avoids \\
Graph entry ($C^{1,\gamma}$ normal graph) & uses & avoids \\
Schauder bootstrap (Lieberman) & uses & avoids \\
Gagliardo--Nirenberg interpolation & uses & avoids \\
Nearly-spherical closure (BDV Thm.~3.3) & uses & avoids \\
Fuglede $H^{1/2}$ expansion & uses & avoids \\
Free-boundary regularity (Alt--Caffarelli) & uses & avoids \\\hline
BV coarea on $\partial^*E_t$ for a.e.\ $t$ & not used & uses \\
AGS metric framework, $|\dot F_\rho|_\X$ & not used & uses \\
Centroid kinematic identity (no homothety) & not used & uses \\
Space--time $\Pi$ identity & not used & uses \\
Slicing-then-squaring (good rays / bad rays) & not used & uses \\
Slice-rearrangement $|E\Delta B|^2\le C\Pi$ & not used & uses \\
Sobolev--Sard (Figalli--Maggi App.~A \cite{FigalliMaggi2011}) & not used & uses \\\hline
BDV bounded reduction \cite[\S5]{BraDeP2017} & uses & uses \\
Kohler--Jobin transfer & uses & uses \\\hline
\end{tabular}
\end{center}

The methodological gain of Plan~2 is that no PDE regularity beyond
$W^{2,p}_{\rm loc}$, no free-boundary regularity, no selection,
no graph parametrisation, and no nearly-spherical closure is used.
The price is the new infrastructure in Blocks~(M)--($\Pi$)--(K).

\section{Open issues and audit trail}\label{sec:audit}

The Plan~2 chain has been audited in three waves on the markdown
source side; each post-audit correction is reflected in the
\texttt{WIP\_X.tex} files.

\textbf{Wave~1}\footnote{Internal audit trail: source file
\texttt{agent6-adversarial-audit.md}.} flagged five issues: centre-choice ambiguity (Fraenkel vs.\ centroid); reliance
on the open pointwise per-$\rho$ radial trace (R); missing
Sobolev--Sard citation; broken Cauchy--Schwarz in (G4)~§5.2;
notational drift between centre choices.

\textbf{Wave~2}\footnote{Internal audit trail: source files
\texttt{wave2-*}.} addressed: the centroid choice
$\zbar=|E_\rho|^{-1}\int_{E_\rho}x$; the integrated form (B$^2$-int)
replacing pointwise (R); the Figalli--Maggi App.~A Sobolev--Sard
citation; the kinetic bound (K) rederivation from (B$^2$-int).

\textbf{Wave~3}\footnote{Internal audit trail: source files
\texttt{wave3-F,G,H,J}.} introduced two new structural inputs: (F) the $H^1$-in-$\rho$ centroid bound
(\textsc{CentroidBound}); and the slicing-then-squaring
decomposition (\textsc{SpaceTimeIdentity} §5). The latter replaces
the broken Cauchy--Schwarz of (G4) §5.2 by the linear good-ray
equality and the slice-rearrangement; squaring before integrating
preserves the sharp $\delta$ rate.

\textbf{Status.} The Plan~2 proof closes in its integrated form
only. Pointwise per-$\rho$ statements (R) and (B$^2$) are not
invoked. The integrated form, controlled by (F) of
\textsc{CentroidBound} and (B$^2$-int) of \textsc{SpaceTimeIdentity},
suffices for the metric-trace and boundary-layer-transfer chain.
This is the unconditional honest content of Plan~2.

\begin{thebibliography}{99}

\bibitem{LevelSetIdentity}
Plan~2 building block \textsc{LevelSetIdentity},
\texttt{WIP\_LevelSetIdentity.tex}.

\bibitem{MetricFramework}
Plan~2 building block \textsc{MetricFramework},
\texttt{WIP\_MetricFramework.tex}.

\bibitem{CentroidBound}
Plan~2 building block \textsc{CentroidBound},
\texttt{WIP\_CentroidBound.tex}.

\bibitem{SpaceTimeIdentity}
Plan~2 building block \textsc{SpaceTimeIdentity},
\texttt{WIP\_SpaceTimeIdentity.tex}.

\bibitem{WeightedMetricTrace}
Plan~2 building block \textsc{WeightedMetricTrace},
\texttt{WIP\_WeightedMetricTrace.tex}.

\bibitem{BoundaryLayerTransfer}
Plan~2 building block \textsc{BoundaryLayerTransfer},
\texttt{WIP\_BoundaryLayerTransfer.tex}.

\bibitem{GlobalAssembly}
Plan~2 building block \textsc{GlobalAssembly},
\texttt{WIP\_GlobalAssembly.tex}.

\bibitem{Plan1Master}
Plan~1 master document (companion route), \texttt{Final/master.tex}.

\bibitem{BDV-BoundedReduction-Plan1}
Plan~1 building block \textsc{BoundedReduction},
\texttt{Final/BoundedReduction.tex}.

\bibitem{KohlerJobinTransfer-Plan1}
Plan~1 building block \textsc{KohlerJobinTransfer},
\texttt{Final/KohlerJobinTransfer.tex}.

\bibitem{BraDeP2017} L.~Brasco, G.~De~Philippis, B.~Velichkov,
\emph{Faber--Krahn inequalities in sharp quantitative form},
Duke Math.\ J.\ \textbf{164} (2015), 1777--1831.

\bibitem{FMP2008} N.~Fusco, F.~Maggi, A.~Pratelli,
\emph{The sharp quantitative isoperimetric inequality},
Ann.\ of Math.\ \textbf{168} (2008), 941--980.

\bibitem{FMPLaplace} N.~Fusco, F.~Maggi, A.~Pratelli,
Ann.\ Sc.\ Norm.\ Super.\ Pisa Cl.\ Sci.\ (5) \textbf{8} (2009), 51--71.

\bibitem{FJ2014} N.~Fusco, V.~Julin, Calc.\ Var.\ PDE \textbf{50} (2014), 925--937.

\bibitem{CL2012} M.~Cicalese, G.~P.~Leonardi,
Arch.\ Ration.\ Mech.\ Anal.\ \textbf{206} (2012), 617--643.

\bibitem{AFM2013} E.~Acerbi, N.~Fusco, M.~Morini,
Comm.\ Math.\ Phys.\ \textbf{322} (2013), 515--557.

\bibitem{AGS2008} L.~Ambrosio, N.~Gigli, G.~Savar\'e,
\emph{Gradient Flows in Metric Spaces…}, Birkh\"auser, 2008.

\bibitem{Brasco2014} L.~Brasco,
ESAIM Control Optim.\ Calc.\ Var.\ \textbf{20} (2014), 315--338.

\bibitem{KohlerJobin} M.-T.~Kohler-Jobin,
Z.\ Angew.\ Math.\ Phys.\ \textbf{29} (1978), 757--766.

\bibitem{Talenti1976} G.~Talenti,
Ann.\ Mat.\ Pura Appl.\ \textbf{110} (1976), 353--372.

\bibitem{Maggi2012} F.~Maggi, Cambridge Univ.\ Press, 2012.

\bibitem{AFP2000} L.~Ambrosio, N.~Fusco, D.~Pallara, Oxford, 2000.

\bibitem{DeGiorgi1958} E.~De~Giorgi,
Atti Accad.\ Naz.\ Lincei \textbf{5} (1958), 33--44.

\bibitem{FigalliMaggi2011} A.~Figalli, F.~Maggi,
Arch.\ Ration.\ Mech.\ Anal.\ \textbf{201} (2011), 143--207.

\bibitem{FMP2010} A.~Figalli, F.~Maggi, A.~Pratelli,
Invent.\ Math.\ \textbf{182} (2010), 167--211.

\bibitem{Faber1923} G.~Faber, Bayer.\ Akad.\ Wiss.\ M\"unchen 1923, 169--172.

\bibitem{Krahn1925} E.~Krahn, Math.\ Ann.\ \textbf{94} (1925), 97--100.

\bibitem{KawohlBook} B.~Kawohl, Springer LNM 1150, 1985.

\bibitem{HenrotBook} A.~Henrot, Birkh\"auser, 2006.

\bibitem{BucurButtazzo} D.~Bucur, G.~Buttazzo, Birkh\"auser, 2005.

\bibitem{HansenNadirashvili1994} W.~Hansen, N.~Nadirashvili,
Potential Anal.\ \textbf{3} (1994), 1--14.

\bibitem{Melas1992} A.~D.~Melas,
J.\ Differential Geom.\ \textbf{36} (1992), 19--33.

\bibitem{Bhattacharya2001} T.~Bhattacharya,
Electron.\ J.\ Differential Equations \textbf{35} (2001), 1--15.

\bibitem{AKN} M.~Allen, D.~Kriventsov, R.~Neumayer,
Ars Inveniendi Analytica (2023).

\bibitem{AC1981} H.~W.~Alt, L.~A.~Caffarelli,
J.\ Reine Angew.\ Math.\ \textbf{325} (1981), 105--144.

\bibitem{Fuglede} B.~Fuglede,
Trans.\ AMS \textbf{314} (1989), 619--638.

\bibitem{KN} D.~Kinderlehrer, L.~Nirenberg,
Ann.\ Scuola Norm.\ Sup.\ Pisa (4) \textbf{4} (1977), 373--391.

\bibitem{Li86} G.~M.~Lieberman,
J.\ Math.\ Anal.\ Appl.\ \textbf{113} (1986), 422--440.

\end{thebibliography}

\end{document}
